\begin{activity}[Warm-up]{Wedge vs.\ Orogen}

\section*{Purpose}
Both accretionary wedges and orogens form at convergent margins, but they experience very different thermal histories. Before we build a detailed model for one of them later today, predict how surface heat flow evolves over the life of each.

\section*{Accretionary wedge}
Sediment and crust are continuously scraped off a subducting oceanic plate and stacked at the trench. The wedge is built from material that has spent little time near the surface; it arrives cold, wet, and is rapidly buried by whatever is scraped off next.

\begin{itemize}
  \item While subduction is active, is heat flow through the wedge likely to be higher or lower than a stable craton? What is actually supplying heat to this material, and what is working against it staying warm?
  \answerbox[2cm]

  \item Once subduction stops---no more cold material arriving, no more rapid burial---what happens to heat flow? Does it recover to the same level as before subduction started, or could it overshoot?
  \answerbox[2cm]
\end{itemize}

\section*{Orogen}
Continental collision thickens the crust; doubling it is not unusual. The thickened crust contains more total radiogenic material than before, buried deeper than it used to be.

\begin{itemize}
  \item Right after thickening, is the crust hotter or colder than it will eventually become? \emph{Hint:} think about how long it takes buried heat production to actually raise temperatures at depth.
  \answerbox[2cm]

  \item Once the orogen has had a long time to adjust, how does its eventual heat flow compare to a stable craton: higher, lower, or about the same?
  \answerbox[2cm]
\end{itemize}

\section*{Sketch both}
On the axes below, sketch your predicted heat flow through time for \emph{both} settings, on the same plot. The dashed line marks typical cratonic heat flow, for reference.

\begin{center}
\begin{tikzpicture}[x=1cm, y=1cm]
  \draw[thick,->] (0,0) -- (9.5,0) node[right] {time};
  \draw[thick,->] (0,0) -- (0,6.5) node[above] {heat flow};
  \draw[dashed, gray] (0,2) -- (9,2);
  \node[gray, right, font=\scriptsize] at (9,2) {craton};
\end{tikzpicture}
\end{center}

\noindent\textit{We will not resolve the wedge case in detail today -- keep your prediction in mind as a point of contrast. The orogen case is looked at in greater detail in scenarios examined later.}
\end{activity}